\documentclass{article}

\usepackage{booktabs}
\usepackage{array}
\usepackage{amsfonts}
\usepackage{amsmath}
\usepackage{graphicx}
\usepackage{hyperref}
\usepackage{url}
\usepackage[margin=1in]{geometry}
\usepackage[numbers,compress,sort]{natbib}

\title{Validating Causal Steering in Protein Language Models: \\
Proxy Accuracy Does Not Guarantee Steering Success}

\author{Dmitriy Ivkov \\
University of Michigan, Ann Arbor \\
\texttt{divkov@umich.edu}}
\date{}

\begin{document}
\maketitle

\begin{abstract}
Activation steering changes model output by modifying internal
activations. Protein steering studies often use scoring proxies because
generated sequences cannot be tested in a wet-lab assay during every
model evaluation. We ask whether agreement between a proxy and
experimental labels also indicates that the proxy will provide a useful
steering target. We test five protein properties in ESM2-650M with
difference-of-means steering vectors and a fixed six-requirement
evaluation protocol. The binding-affinity proxy has the strongest
held-out correlation in the study ($r=0.80$), but its steering direction
has no measurable effect. A weaker catalytic-activity proxy ($r=0.22$)
produces an effect that repeats across three random seeds. We find a
similar mismatch for individual attention heads: strong correlation
with known 3D contacts does not imply a large effect under ablation. Two
additional checks change how we interpret other results. A
residue-exclusion check gives different conclusions across seeds, and an
apparent expression-yield effect uses a direction that overlaps with the
intrinsic-disorder direction. These results show that proxy validation
tests the score, while direct intervention tests are needed to evaluate
the steering direction.
\end{abstract}

\section{Introduction}

Difference-of-means activation steering modifies a model's residual
stream during generation. The steering direction is obtained by
subtracting the mean activation of a low-property group from the mean
activation of a high-property group. This method has reproduced a known
thermostability effect in ESM2-650M \citep{huang2025steering}.

That result leaves an important question open when the method is applied
to a new protein property. Generated proteins cannot be evaluated with a
wet-lab assay during every model run, so steering studies often depend
on a scoring proxy. A strong correlation between the proxy and
experimental labels shows that the score contains information about the
property. However, it does not show that a difference-of-means vector
built from that property will change the model in the intended way.

We test five properties under the same protocol, and four of them reach
the steering stage. The order of their proxy correlations does not match
the order of their steering effects. The binding-affinity proxy has the
strongest held-out correlation, but its steering vector has no
detectable effect. The catalytic-activity proxy has a much weaker
correlation, yet its effect is consistent across intervention strengths
and random seeds.

We test a related assumption at the level of individual model components.
An attention head that correlates strongly with known protein contacts
is not unusually important when it is ablated. Together, these
experiments distinguish evidence that a model represents a property from
evidence that a particular intervention controls that property.

This paper evaluates the reliability of steering evidence rather than
introducing a new steering algorithm. We apply one protocol to five
properties, investigate why two results fail or remain inconclusive, and
document two analysis choices that would have changed our conclusions.

\section{Related Work}

\citet{huang2025steering} introduced difference-of-means activation
steering for protein language models and validated the method on
thermostability. We apply the same intervention to five new properties
and test whether proxy accuracy predicts steering success.

\citet{vig2021bertology} identified attention heads whose attention
patterns correlate with 3D protein contacts. The original study
described this evidence as associative. We reproduce their analysis,
test the most enriched head in our data by causal ablation, and extend
the analysis to all 480 heads in the model.

\citet{adams2025mechanistic} used linear probes to identify
thermostability and localization information in sparse-autoencoder
features from ESM-2. Their steering experiment tested whether a selected
latent produced more coherent sequence changes than a random latent. It
did not compare the result with an experimental property label. Our
study addresses a related but different question: whether correlational
evidence predicts the result of a downstream intervention.

\section{Methods}
\label{sec:method}

\subsection{Study design}

We study binding affinity, catalytic activity, intrinsic disorder,
immunogenicity, and expression yield. We evaluate each property through
the same sequence of steps. First, we validate a property-specific score
against experimental labels. If the score passes that check, we
construct a steering vector, generate sequences at several intervention
strengths, and compare the learned vector with a matched random control.
The immunogenicity study stops after proxy validation because none of
its candidate scores generalizes to the experimental endpoint.

\subsection{Activation steering}

For each target property, we divide real protein sequences into
high-property and low-property groups. At every layer of ESM2-650M, we
compute the mean residual-stream activation for each group and subtract
the two means. The resulting vector is added to the residual stream
during sequence generation. A scalar $\alpha$ controls the intervention
strength.

Each generated sequence is evaluated with the property-specific
compositional proxy. We compare the property-derived steering direction
with a random direction of the same norm. This comparison tests whether
the measured change depends on the learned direction rather than only on
the size of the intervention.

\subsection{Six evaluation requirements}

We fix six requirements before testing the new properties:

\begin{enumerate}
\item The property-derived direction must outperform a random direction
  of the same norm in a paired bootstrap test with 10,000 resamples.
\item The effect must increase monotonically across at least three
  intervention strengths in a range where generation remains stable.
\item The effect must remain after the most frequently substituted
  residues are excluded from the score.
\item The scoring proxy must agree with experimental labels on held-out
  data before steering is attempted.
\item The method must outperform an existing baseline when a suitable
  baseline is available.
\item An evaluation of a new property must contain at least 150 held-out
  sequences.
\end{enumerate}

The third requirement is a check for compositional collapse. A steering
direction can appear successful if it repeatedly inserts one residue
that directly increases the proxy score. Excluding the dominant
substitutions tests whether the effect extends beyond that narrow
change.

We use three assessment outcomes. If proxy validation fails, we reject
the proxy and do not attempt steering. If the main steering test
succeeds but the residue-exclusion result fails or changes across seeds,
we classify the evidence as inconclusive. We call a result supported
only when it satisfies every applicable requirement.

\subsection{Attention-head ablation}

We also examine whether correlational evidence predicts the importance
of individual attention heads. Following \citet{vig2021bertology}, we
measure contact enrichment in their protein language model and select
the most enriched head in our data. We ablate that head and compare the
change in masked-residue prediction with a low-correlation control head.
We then repeat the ablation for all 480 heads in the model and rank them
by measured causal effect.

\section{Results}
\label{sec:results}

\subsection{Results across five target properties}

Table~\ref{tab:sweep} summarizes the five target properties. Binding
affinity has the strongest proxy correlation but no measurable steering
effect. Catalytic activity has the weakest proxy correlation among the
four properties that reach the steering stage, yet its effect is larger
and repeats across seeds.

\begin{table}[t]
\centering
\caption{Results for the five target properties. Proxy $r$ is the
Pearson correlation with experimental labels on held-out data, measured
before steering.}
\label{tab:sweep}
\small
\begin{tabular}{@{}>{\raggedright\arraybackslash}p{2.8cm}
                    >{\raggedright\arraybackslash}p{3.2cm}
                    >{\raggedright\arraybackslash}p{1.4cm}
                    >{\raggedright\arraybackslash}p{4.3cm}
                    >{\raggedright\arraybackslash}p{2.8cm}@{}}
\toprule
Target & Dataset & Proxy $r$ & Steering result & Assessment \\
\midrule
Binding affinity & ProteinGym DMS (KRAS/DARPin) & 0.80 & no measurable effect & Unsupported \\
Catalytic activity & DLKcat (cross-protein) & 0.22 & repeats across 3 seeds & Supported \\
Intrinsic disorder & DisProt & 0.45 & main effect 3/3; check 2/3 & Inconclusive \\
Immunogenicity & IEDB (4-tier) & $\leq0.10$ & steering not attempted & Proxy rejected \\
Expression yield & eSol & 0.31 & fails residue-exclusion check & Inconclusive \\
\bottomrule
\end{tabular}
\end{table}

\begin{figure}[t]
\centering
\includegraphics[width=0.72\linewidth]{figures/fig2_proxy_vs_effect.pdf}
\caption{Proxy correlation and steering effect size for the four targets
that reach the steering stage. Effect size is Cohen's $d$ at
$\alpha=0.5$ for the learned direction relative to the matched random
direction. The four points show no positive relation between the two
measurements.}
\label{fig:proxy-vs-effect}
\end{figure}

\subsection{Binding affinity}

We use a ProteinGym deep-mutational-scanning assay
\citep{notin2023proteingym} of KRAS binding to a DARPin. The proxy is
built from position-specific mutational-sensitivity weights learned on
the training split. It scores a sequence according to whether the
wild-type residue is preserved at positions where mutations have a large
measured effect.

The proxy reaches $r=0.80$ on held-out data. A weight-shuffle control
gives $r=-0.001$, which indicates that the position-specific weights
carry the signal. The proxy also generalizes to positions excluded from
weight fitting ($r=0.54$) and from single to double mutants
($r=0.70$).

Despite this validation, the steering direction has no detectable effect
at any tested strength. Generation remains stable, and no strength
produces degenerate sequences. Output collapse therefore does not
explain the result.

\subsection{Catalytic activity}

For catalytic activity, we use a cross-protein enzyme turnover dataset
from DLKcat \citep{li2022dlkcat}. The proxy measures the fraction of
glycine minus the fraction of arginine.

The proxy reaches held-out $r=0.22$, the weakest result among the four
targets that reach the steering stage. Even so, the steering effect is
significant at every tested strength in the stable range. It increases
monotonically, remains after residue exclusion, and repeats across three
random seeds.

\subsection{Intrinsic disorder}
\label{sec:disorder}

We use DisProt annotations \citep{quaglia2022disprot}. The
intrinsic-disorder proxy uses the published TOP-IDP amino-acid scale
\citep{campen2008topidp}. It reaches held-out $r=0.45$. A sliding-window
comparison with per-residue disorder labels gives AUC $0.71$.

The main steering effect is significant and monotonic on all three
random seeds. The residue-exclusion check is less stable. It passes on
two seeds, retaining 36--42\% of the original effect, but fails on the
third, where it retains 10\% and its confidence interval crosses zero.
The effect is repeatable, but the evidence against a simple
compositional artifact is not equally strong on every seed.

\begin{figure}[t]
\centering
\includegraphics[width=\linewidth]{figures/fig1_dose_response.pdf}
\caption{Dose-response curves for catalytic activity and intrinsic
disorder. Both effects increase monotonically across the tested
intervention strengths. Error bars show 95\% bootstrap confidence
intervals.}
\label{fig:dose-response}
\end{figure}

\begin{figure}[t]
\centering
\includegraphics[width=0.75\linewidth]{figures/fig3_seed_robustness.pdf}
\caption{Intrinsic-disorder effects before and after residue exclusion
for three random seeds. The main effect is stable, but the confidence
interval for the residue-excluded score crosses zero on seed 1.}
\label{fig:seed-robustness}
\end{figure}

\subsection{Immunogenicity}

For immunogenicity, we use records from the Immune Epitope Database
\citep{vita2019iedb} and evaluate candidate proxies at four levels:
peptide MHC-II binding affinity, mass-spectrometry-confirmed peptide
presentation, measured T-cell response, and full-length source antigens.
Candidate proxies correlate with the surrogate binding-affinity endpoint
at $r=0.37$--$0.43$ but fall to $r\leq0.10$ on the measured
T-cell-response endpoint.

One proxy for full-length antigens initially reaches $r=0.38$. Further
testing shows that organism identity explains 58\% of this signal. When
organism identity is held out during cross-validation, the correlation
changes to $-0.32$. Because no proxy passes the validation requirement,
we do not run a steering experiment.

\subsection{Expression yield}
\label{sec:expression}

For expression yield, we use the eSol dataset
\citep{niwa2009esol}. The proxy is based on the absolute deviation of net
charge from neutrality, following a published recombinant-solubility
discriminant \citep{wilkinson1991solubility}. It reaches held-out
$r=0.31$--$0.34$ and produces a significant, monotonic steering effect.

The effect disappears when the dominant residue substitutions are
excluded from scoring. One of those residues is part of the proxy
definition. The result is therefore inconclusive under the fixed
protocol.

\subsection{Attention-head ablation}

In our replication of \citet{vig2021bertology}, the most
contact-enriched head has 12.9-fold enrichment over background. Ablating
it produces a change in masked-residue prediction that is statistically
indistinguishable from ablating a low-correlation control head.

The full sweep supports the same interpretation, but it does not provide
a precise rank for this head. The evaluation uses 104 sampled positions,
so many heads have identical measured effects. The selected head shares
its value with 131 other heads, forming a tied block that spans ranks
286--417. At the same resolution, 166 heads show no measurable effect.
Strong contact correlation therefore does not imply a large or unique
contribution to masked-residue prediction.

\section{Failure Analysis and Reliability Checks}
\label{sec:failure}

\subsection{The binding-affinity groups differ at few residues}

The binding-affinity sequences are point mutants of one 188-residue
reference protein. The high- and low-binding groups differ by only one
or two residues per sequence. Their mean per-residue compositional
distance is $0.003$, compared with $0.02$--$0.05$ for the three
cross-protein datasets.

This difference offers a possible explanation for the failed
intervention. The proxy uses position-specific weights, so it can predict
binding labels from a small number of mutations. In contrast, the
steering direction averages activations over groups of nearly identical
sequences. The few mutations that distinguish the groups may contribute
too little to that average. Our evidence supports this explanation as a
possibility, not as a demonstrated mechanism.

\subsection{The expression-yield direction overlaps with disorder}

We compare the expression-yield steering vector with every other target
vector. Its cosine similarity with the intrinsic-disorder direction is
$+0.38$ across all 33 layers and rises to $+0.56$--$0.67$ in the three
deepest layers. This overlap shows that the two directions are not
orthogonal and may capture some of the same compositional variation. It
does not show that disorder causes the expression-yield result. Because
the expression-yield result also fails residue exclusion, we do not count
it as separate evidence for a new steering effect. The comparison
requires one additional embedding pass and is inexpensive relative to a
complete steering sweep.

\subsection{Unstable generation made five-layer steering look stronger}
\label{sec:unstable}

In a follow-up thermostability experiment, we compared steering all 33
layers with steering only five layers previously identified as
important. An early analysis selected the best intervention strength
from the full tested range, including $\alpha\geq1.0$. At the selected
strength, generation under the 33-layer intervention had already
collapsed, while generation under the 5-layer intervention had not.
The 5-layer version appeared better only because it failed at a higher
strength.

We restricted the comparison to a range where both interventions
produced stable sequences and reran the complete experiment. The
5-layer intervention then retained 43--59\% of the full effect. This is
a useful partial result, but it does not support the original claim that
five layers perform as well as all 33.

\subsection{A robustness criterion changed across seeds}

The intrinsic-disorder effect repeats across all three seeds, but its
residue-exclusion check passes on only two. A single run could therefore
support either a pass or a failure for that criterion. Repeating only
the main effect is not enough when a separate artifact check is used to
justify the result.

\section{Discussion}

These results show why representation and intervention should be
evaluated separately. A proxy can accurately summarize experimental
labels without defining a direction that changes model output in the
intended way. Likewise, an attention head can align with a biological
pattern without making a large, individually measurable contribution to
the model's prediction.

Proxy validation remains necessary because it rules out scores that do
not track the target, as the immunogenicity experiment shows. The problem
arises when evidence about the score is also treated as evidence about
the steering vector. The binding-affinity experiment shows that the score
and the vector can depend on different sources of signal.

The experiments suggest several practical checks. Before constructing a
steering vector, researchers should measure how much the high- and
low-property groups differ. Comparisons between interventions should use
the same range of strengths, and every intervention should produce
stable sequences throughout that range. Important results should be
repeated across seeds, including the artifact checks used to support
them. When several properties are studied together, similarities between
their steering vectors can reveal that two apparent effects may not be
independent.

Every successful result in this study is causal only with respect to the
model's activation space. The intervention changes generated sequences
and outperforms a matched random direction, but no generated protein was
tested in a biological assay. A successful protocol result means that
the intervention changes a validated proxy. It does not establish that
the generated protein has the intended biological property.

\section{Limitations}

All experiments use ESM2-650M at one model scale. The relationship
between proxy accuracy and steering effect could differ in other models
or at other scales. The proxies are compositional heuristics rather than
learned predictors. A stronger non-compositional proxy could behave
differently.

We selected each final proxy from several candidates using the same
held-out split on which we report its correlation. The reported
correlations are therefore likely optimistic and should be treated as
upper bounds. This issue affects the absolute values, although it does
not reverse the ordering between the binding-affinity and
catalytic-activity proxies.

We ran multiple significance tests without a formal correction.
Agreement across intervention strengths and seeds reduces this risk but
does not remove it. Seed robustness was checked for only two targets.
We also did not evaluate generated sequences with a structural
plausibility model or a wet-lab assay.

\section{Conclusion}

Across these experiments, higher proxy accuracy does not guarantee more
successful steering. The binding-affinity result provides the clearest
counterexample, and the attention-head ablation shows a related mismatch
at the level of a model component. Correlational evidence can show that a
model represents a property, but reliable steering claims still require
direct tests of the intervention, its artifacts, and its stability.

\section*{LLM usage disclosure}
Claude (Anthropic) assisted with code writing, experiment scripts,
literature searches, and manuscript drafting. The author chose the study
design, ran and checked the experiments, and verified each numerical
claim against committed outputs.

\section*{Code and data availability}
Code, data-processing scripts, saved results, and instructions for
reproducing the experiments are available at
\url{https://github.com/DIvkov575/esmplm-steering}.

\bibliographystyle{plainnat}
\bibliography{reference}

\end{document}
